% Options for packages loaded elsewhere
\PassOptionsToPackage{unicode}{hyperref}
\PassOptionsToPackage{hyphens}{url}
\PassOptionsToPackage{dvipsnames,svgnames,x11names}{xcolor}
%
\documentclass[
  12pt]{article}

\usepackage{amsmath,amssymb}
\usepackage{iftex}
\ifPDFTeX
  \usepackage[T1]{fontenc}
  \usepackage[utf8]{inputenc}
  \usepackage{textcomp} % provide euro and other symbols
\else % if luatex or xetex
  \usepackage{unicode-math}
  \defaultfontfeatures{Scale=MatchLowercase}
  \defaultfontfeatures[\rmfamily]{Ligatures=TeX,Scale=1}
\fi
\usepackage{lmodern}
\ifPDFTeX\else  
    % xetex/luatex font selection
\fi
% Use upquote if available, for straight quotes in verbatim environments
\IfFileExists{upquote.sty}{\usepackage{upquote}}{}
\IfFileExists{microtype.sty}{% use microtype if available
  \usepackage[]{microtype}
  \UseMicrotypeSet[protrusion]{basicmath} % disable protrusion for tt fonts
}{}
\makeatletter
\@ifundefined{KOMAClassName}{% if non-KOMA class
  \IfFileExists{parskip.sty}{%
    \usepackage{parskip}
  }{% else
    \setlength{\parindent}{0pt}
    \setlength{\parskip}{6pt plus 2pt minus 1pt}}
}{% if KOMA class
  \KOMAoptions{parskip=half}}
\makeatother
\usepackage{xcolor}
\setlength{\emergencystretch}{3em} % prevent overfull lines
\setcounter{secnumdepth}{5}
% Make \paragraph and \subparagraph free-standing
\makeatletter
\ifx\paragraph\undefined\else
  \let\oldparagraph\paragraph
  \renewcommand{\paragraph}{
    \@ifstar
      \xxxParagraphStar
      \xxxParagraphNoStar
  }
  \newcommand{\xxxParagraphStar}[1]{\oldparagraph*{#1}\mbox{}}
  \newcommand{\xxxParagraphNoStar}[1]{\oldparagraph{#1}\mbox{}}
\fi
\ifx\subparagraph\undefined\else
  \let\oldsubparagraph\subparagraph
  \renewcommand{\subparagraph}{
    \@ifstar
      \xxxSubParagraphStar
      \xxxSubParagraphNoStar
  }
  \newcommand{\xxxSubParagraphStar}[1]{\oldsubparagraph*{#1}\mbox{}}
  \newcommand{\xxxSubParagraphNoStar}[1]{\oldsubparagraph{#1}\mbox{}}
\fi
\makeatother

\usepackage{color}
\usepackage{fancyvrb}
\newcommand{\VerbBar}{|}
\newcommand{\VERB}{\Verb[commandchars=\\\{\}]}
\DefineVerbatimEnvironment{Highlighting}{Verbatim}{commandchars=\\\{\}}
% Add ',fontsize=\small' for more characters per line
\usepackage{framed}
\definecolor{shadecolor}{RGB}{241,243,245}
\newenvironment{Shaded}{\begin{snugshade}}{\end{snugshade}}
\newcommand{\AlertTok}[1]{\textcolor[rgb]{0.68,0.00,0.00}{#1}}
\newcommand{\AnnotationTok}[1]{\textcolor[rgb]{0.37,0.37,0.37}{#1}}
\newcommand{\AttributeTok}[1]{\textcolor[rgb]{0.40,0.45,0.13}{#1}}
\newcommand{\BaseNTok}[1]{\textcolor[rgb]{0.68,0.00,0.00}{#1}}
\newcommand{\BuiltInTok}[1]{\textcolor[rgb]{0.00,0.23,0.31}{#1}}
\newcommand{\CharTok}[1]{\textcolor[rgb]{0.13,0.47,0.30}{#1}}
\newcommand{\CommentTok}[1]{\textcolor[rgb]{0.37,0.37,0.37}{#1}}
\newcommand{\CommentVarTok}[1]{\textcolor[rgb]{0.37,0.37,0.37}{\textit{#1}}}
\newcommand{\ConstantTok}[1]{\textcolor[rgb]{0.56,0.35,0.01}{#1}}
\newcommand{\ControlFlowTok}[1]{\textcolor[rgb]{0.00,0.23,0.31}{\textbf{#1}}}
\newcommand{\DataTypeTok}[1]{\textcolor[rgb]{0.68,0.00,0.00}{#1}}
\newcommand{\DecValTok}[1]{\textcolor[rgb]{0.68,0.00,0.00}{#1}}
\newcommand{\DocumentationTok}[1]{\textcolor[rgb]{0.37,0.37,0.37}{\textit{#1}}}
\newcommand{\ErrorTok}[1]{\textcolor[rgb]{0.68,0.00,0.00}{#1}}
\newcommand{\ExtensionTok}[1]{\textcolor[rgb]{0.00,0.23,0.31}{#1}}
\newcommand{\FloatTok}[1]{\textcolor[rgb]{0.68,0.00,0.00}{#1}}
\newcommand{\FunctionTok}[1]{\textcolor[rgb]{0.28,0.35,0.67}{#1}}
\newcommand{\ImportTok}[1]{\textcolor[rgb]{0.00,0.46,0.62}{#1}}
\newcommand{\InformationTok}[1]{\textcolor[rgb]{0.37,0.37,0.37}{#1}}
\newcommand{\KeywordTok}[1]{\textcolor[rgb]{0.00,0.23,0.31}{\textbf{#1}}}
\newcommand{\NormalTok}[1]{\textcolor[rgb]{0.00,0.23,0.31}{#1}}
\newcommand{\OperatorTok}[1]{\textcolor[rgb]{0.37,0.37,0.37}{#1}}
\newcommand{\OtherTok}[1]{\textcolor[rgb]{0.00,0.23,0.31}{#1}}
\newcommand{\PreprocessorTok}[1]{\textcolor[rgb]{0.68,0.00,0.00}{#1}}
\newcommand{\RegionMarkerTok}[1]{\textcolor[rgb]{0.00,0.23,0.31}{#1}}
\newcommand{\SpecialCharTok}[1]{\textcolor[rgb]{0.37,0.37,0.37}{#1}}
\newcommand{\SpecialStringTok}[1]{\textcolor[rgb]{0.13,0.47,0.30}{#1}}
\newcommand{\StringTok}[1]{\textcolor[rgb]{0.13,0.47,0.30}{#1}}
\newcommand{\VariableTok}[1]{\textcolor[rgb]{0.07,0.07,0.07}{#1}}
\newcommand{\VerbatimStringTok}[1]{\textcolor[rgb]{0.13,0.47,0.30}{#1}}
\newcommand{\WarningTok}[1]{\textcolor[rgb]{0.37,0.37,0.37}{\textit{#1}}}

\providecommand{\tightlist}{%
  \setlength{\itemsep}{0pt}\setlength{\parskip}{0pt}}\usepackage{longtable,booktabs,array}
\usepackage{calc} % for calculating minipage widths
% Correct order of tables after \paragraph or \subparagraph
\usepackage{etoolbox}
\makeatletter
\patchcmd\longtable{\par}{\if@noskipsec\mbox{}\fi\par}{}{}
\makeatother
% Allow footnotes in longtable head/foot
\IfFileExists{footnotehyper.sty}{\usepackage{footnotehyper}}{\usepackage{footnote}}
\makesavenoteenv{longtable}
\usepackage{graphicx}
\makeatletter
\newsavebox\pandoc@box
\newcommand*\pandocbounded[1]{% scales image to fit in text height/width
  \sbox\pandoc@box{#1}%
  \Gscale@div\@tempa{\textheight}{\dimexpr\ht\pandoc@box+\dp\pandoc@box\relax}%
  \Gscale@div\@tempb{\linewidth}{\wd\pandoc@box}%
  \ifdim\@tempb\p@<\@tempa\p@\let\@tempa\@tempb\fi% select the smaller of both
  \ifdim\@tempa\p@<\p@\scalebox{\@tempa}{\usebox\pandoc@box}%
  \else\usebox{\pandoc@box}%
  \fi%
}
% Set default figure placement to htbp
\def\fps@figure{htbp}
\makeatother

\addtolength{\oddsidemargin}{-.5in}%
\addtolength{\evensidemargin}{-1in}%
\addtolength{\textwidth}{1in}%
\addtolength{\textheight}{1.7in}%
\addtolength{\topmargin}{-1in}%
\usepackage{url}
\usepackage{xurl}
\usepackage{booktabs}
\usepackage{tabularx}
\usepackage{xltabular}
\usepackage{array}
\usepackage{float}
\keepXColumns
\usepackage{ragged2e}
\newcolumntype{T}[1]{>{\RaggedRight\ttfamily\arraybackslash}p{#1}}
\newcommand{\us}{\_\allowbreak}
\makeatletter
\@ifpackageloaded{caption}{}{\usepackage{caption}}
\AtBeginDocument{%
\ifdefined\contentsname
  \renewcommand*\contentsname{Table of contents}
\else
  \newcommand\contentsname{Table of contents}
\fi
\ifdefined\listfigurename
  \renewcommand*\listfigurename{List of Figures}
\else
  \newcommand\listfigurename{List of Figures}
\fi
\ifdefined\listtablename
  \renewcommand*\listtablename{List of Tables}
\else
  \newcommand\listtablename{List of Tables}
\fi
\ifdefined\figurename
  \renewcommand*\figurename{Figure}
\else
  \newcommand\figurename{Figure}
\fi
\ifdefined\tablename
  \renewcommand*\tablename{Table}
\else
  \newcommand\tablename{Table}
\fi
}
\@ifpackageloaded{float}{}{\usepackage{float}}
\floatstyle{ruled}
\@ifundefined{c@chapter}{\newfloat{codelisting}{h}{lop}}{\newfloat{codelisting}{h}{lop}[chapter]}
\floatname{codelisting}{Listing}
\newcommand*\listoflistings{\listof{codelisting}{List of Listings}}
\makeatother
\makeatletter
\makeatother
\makeatletter
\@ifpackageloaded{caption}{}{\usepackage{caption}}
\@ifpackageloaded{subcaption}{}{\usepackage{subcaption}}
\makeatother

\usepackage[]{natbib}
\bibliographystyle{agsm}
\usepackage{bookmark}

\IfFileExists{xurl.sty}{\usepackage{xurl}}{} % add URL line breaks if available
\urlstyle{same} % disable monospaced font for URLs
\hypersetup{
  pdftitle={Project Title (Placeholder)},
  pdfauthor={Camden Bergquist},
  colorlinks=true,
  linkcolor={blue},
  filecolor={Maroon},
  citecolor={Blue},
  urlcolor={Blue},
  pdfcreator={LaTeX via pandoc}}



\begin{document}


\def\spacingset#1{\renewcommand{\baselinestretch}%
{#1}\small\normalsize} \spacingset{1}


%%%%%%%%%%%%%%%%%%%%%%%%%%%%%%%%%%%%%%%%%%%%%%%%%%%%%%%%%%%%%%%%%%%%%%%%%%%%%%

\date{January 30, 2026}
\title{\bf Project Title (Placeholder)}
\author{
Camden Bergquist\\
}
\maketitle

\bigskip
\bigskip
\begin{abstract}

\end{abstract}


\newpage
\spacingset{1.9} % DON'T change the spacing!

\renewcommand*\contentsname{Table of contents}
{
\hypersetup{linkcolor=}
\setcounter{tocdepth}{3}
\tableofcontents
}
\listoffigures
\listoftables

Make sure every code chunk has echo=FALSE in its preamble unless you
need to highlight the code for some reason in the paper.

\section{Abstract}\label{abstract}

Complete this after all other sections have been completed.

\newpage

\section{Introduction}\label{introduction}

Competitive games have become a major source of large-scale behavioral
data, enabling researchers to study strategic decision-making, skill
development, and population-level trends using observational logs at
scales unheard of compared to controlled studies
\citep{edelmann2020css, watts2014css}. Chess is one of the most famous
competitive games in the world, and the relatively recent growth of
online platforms has created datasets with millions to billions of games
played under standardized rules and well-defined competitive contexts
(time controls, ratings, game outcomes, etc.)
\citep{watts2014css, lichessDatabase}. One of the most widely used
sources is the Lichess Open Database, which provides public monthly
exports of rated games played on lichess.org in PGN format
\citep{lichessDatabase, lichessDatabaseGitHub}. These exports make it
possible to study how strategy and performance vary across skill levels
and time controls in a modern, online environment.

This project intends to use Lichess game data to investigate
contemporary metagame behavior and player development. Here, `metagame'
refers to the population-level distribution of strategic choices and
styles observed in practice -- such as opening families, risk
preferences, and game outcomes -- within a defined time window and
player population. In addition to describing how strategic behavior
differs by rating and time control, the project will attempt to connect
these strategic signatures to measures of player progression over time.
Because chess is a repeatedly played, fully competitive setting with a
mature rating infrastructure, it offers a natural testbed for
quantifying skill and comparing behavior across different players
\citep{lichessRatingSystems, glickmanGlicko2Desc}.

A central component of modern chess analytics is the rating system used
to summarize player strength. The ELO rating system -- a system
originally developed for chess -- models player performance
probabilistically, and remains a foundational reference point for
interpreting skill differences and expected outcomes in competitive
settings \citep{elo1978rating, fideElo2025}. Lichess uses a modified
version of this system called `Glicko-2', which, in addition to rating,
tracks uncertainty via rating deviation and volatility
\citep{lichessRatingSystems, glickmanGlicko2Desc}. Ratings are
especially useful for observational studies because they provide a
consistent scale by which to stratify players and construct longitudinal
measures of improvement \citep{glickmanGlicko2Desc}.

Overall, this project attempts to marry data science, computational
social science, and strategic behavior: It treats online chess as a
real-world environment where strategic choices and performance outcomes
can be measured at scale and subsequently analyzed using statistical
learning methods \citep{edelmann2020css, watts2014css}.

\subsection{Background}\label{background}

A growing research literature uses large online chess datasets to study
human performance, decision-making, and strategic patterns. For example,
Chowdhary et al.~leverage an online chess dataset at the scale of over
100 million games to quantify performance dynamics and career-like
trajectories, illustrating how large game archives support
population-level behavioral measurement across skill levels
\citep{chowdhary2023performance}.

One important line of work focuses on time usage and decision quality
under time constraints. Online chess platforms record both move
sequences and clock information, which allows researchers to study how
players allocate thinking time and how time pressure shapes decisions.
Russek et al.~analyze millions of Lichess games to relate move times to
the estimated value of additional computation, and show that time
allocation patterns vary systematically with playing strength and time
control \citep{russek2025thinking}. These findings support the idea that
online chess is a rich setting for studying bounded rationality and
performance under constraints, especially when comparing bullet, blitz,
rapid, and classical formats \citep{russek2025thinking}.

Chess opening choice is also an active research topic. Because openings
are discrete, labelable strategic choices that occur early in every
game, they are well suited to population-level analysis. De Marzo and
Servedio use online chess community data to quantify opening similarity
and complexity, including representing openings in network form and
studying how opening-related structure can be used to forecast future
opening adoption \citep{demarzo2023openings}. This motivates analyzing
metagame development in terms of shifting opening popularity and success
rates across time, player skill, and time controls
\citep{demarzo2023openings, chowdhary2023performance}.

Finally, there is an emerging interest in identifying player styles or
archetypes from game logs. While traditional chess analysis often
focuses on engine evaluation of positions, a complementary approach is
to extract behavioral features from move sequences and outcomes (e.g.,
game length, opening selection, and other style proxies) and use
statistical learning to characterize player types. For instance,
McIlroy-Young et al.~demonstrate that models trained on move sequences
can capture individual-specific signatures (``style'') in chess
decisions, providing evidence that player behavior is meaningfully
distinguishable from game records alone
\citep{mcilroyyoung2021stylometry}. This creates a bridge between
strategic analysis (what players choose to do) and longitudinal outcomes
(how players improve or persist over time), which aligns closely with
the goals of this project.

Taken together, existing work suggests that large-scale online chess
datasets support three relevant kinds of analysis: (1) descriptive
metagame analysis of strategic choices across populations, (2) modeling
how constraints like time pressure affect behavior and outcomes, and (3)
data-driven characterization of player types and their evolution
\citep{chowdhary2023performance, russek2025thinking, demarzo2023openings, mcilroyyoung2021stylometry}.
This project intends to build on those themes by using Lichess data to
analyze modern opening/style patterns and connect them to measures of
player performance and development over time.

\newpage

\section{The Data}\label{the-data}

\subsection{Source}\label{sec-source}

Describe where this data comes from, why it was collected, who collected
it, and how it was collected. A Citation may be useful here.

Lichess is the largest open-source chess server is the world, and
releases free, monthly exports of their game data under the Creative
Commons CCO license \citep{lichessDatabase}. The site requires a log of
the game data in order to function while a game is being played, but
holds onto it after the fact to allow users to analyze their games after
playing them. Though Lichess offers a powerful API, I elected to
download data directly from their website in \texttt{.pgn.zst} format.
The PGN (portable game notation) format is a plaintext format for chess
games, containing moves, headers like players, date, and event, and even
optional annotation. The \texttt{.zst} file extension (sometimes also
written as \texttt{zstd}) represents Zstandard compression, which is a
modern compressor similar to gzip or xz that's commonly used because
it's fast and compresses well; think of it as a variation on a ZIP or
RAR file. Zstandard compression can be expected to shrink PGN files by 5
times on its faster, less-efficient settings, and even more on it's
slower ones.

For the purposes of this project, I downloaded 24 months of data
spanning from January 2024 to December 2025, amounting to 677 gigabytes
of raw data when compressed. Obviously, nearly 700 gigabytes of
compressed data -- 3\textasciitilde4 terabytes of uncompressed PGN data
-- is an insane amount to try to work with, capstone project or not. I
have a powerful home computer, but that kind of scale requires something
closer to a high-end server than a personal computer; it won't fly.
Instead of tackling the problem head-on, I constructed a hybrid R/Python
pipeline that downloads the data, deterministically samples and parses
it into small chunks (in the form of parquet files), and then builds a
selection of smaller, ready-for-analysis datasets from those chunks.
Individually, no one dataset has the 20+ variables required in the
project outline, but since I've built multiple datasets for multiple
different types of analysis, they collectively meet that benchmark

The download and parsing section of the pipeline handles the following:

\begin{itemize}
\tightlist
\item
  Reads a config file \texttt{pipeline.yml} that controls the scope and
  parameters of the pipeline's operation.
\item
  Downloads monthly Lichess `standard rated' PGN archives for each month
  specified in the YAML.
\item
  Decompresses each Zstandard file as a stream, scanning games
  one-by-one so it doesn't run into RAM bottlenecks.
\item
  Parses game headers first to decide whether a game is relevant before
  doing any heavy work (not all games are actually getting sampled, so
  it skips over the ones not being used).

  \begin{itemize}
  \tightlist
  \item
    Extracts core metadata (players, ratings, time control, result,
    termination, opening, timestamps).
  \item
    Buckets each game into blitz/rapid/classical using the Lichess
    time-control convention.
  \end{itemize}
\item
  Applies filtering rules from the YAML, which in this case are:

  \begin{itemize}
  \tightlist
  \item
    Only include blitz/rapid/classical time controls.
  \item
    Only include games from players in the 800--2600 rating range.
  \item
    Drops unwanted terminations (aborted games).
  \end{itemize}
\item
  Uses deterministic hash-based sampling to create reproducible subsets
  without needing to store the whole population:

  \begin{itemize}
  \tightlist
  \item
    A meta sample keyed by \texttt{game\_id} for population/metagame
    summaries, at a rate of 5/500 games.
  \item
    A player cohort keyed by \texttt{username} for longitudinal/player
    tracking at a rate of 3/500 players.
  \item
    A moves sample keyed by \texttt{game\_id} for clock/time-per-move
    data at a rate of 5/10000 moves.
  \end{itemize}
\item
  Deterministically anonymizes usernames based on a stable, salted hash.

  \begin{itemize}
  \tightlist
  \item
    Cohort selection is computed on the real username before
    anonymization, so it remains reproducible.
  \end{itemize}
\item
  Writes month-partitioned parquet outputs for game-level data.
\end{itemize}

The dataset-building portion of the pipeline takes those outputs and,
using DuckDB to read and write parquet files on the disk, builds seven
datasets from them:

\begin{itemize}
\tightlist
\item
  Player-level longitudinal summary tables from \texttt{games\_cohort}:

  \begin{itemize}
  \tightlist
  \item
    \texttt{player\_month} (\texttt{pm}): player * month, 249,761
    observations.
  \item
    \texttt{player\_month\_tc} (\texttt{pm\_tc}): player * month * time
    control, 320,869 observations.
  \end{itemize}
\item
  Player-opening `repertoire' table from \texttt{games\_cohort}:

  \begin{itemize}
  \tightlist
  \item
    \texttt{player\_opening\_month\_tc} (\texttt{pom\_tc}): player *
    month * time control * opening, 7,167,580 observations.
  \end{itemize}
\item
  Metagame tables from \texttt{games\_meta}:

  \begin{itemize}
  \tightlist
  \item
    \texttt{meta\_month\_tc} (\texttt{mm\_tc}): month * time control, 72
    observations.
  \item
    \texttt{meta\_opening\_month\_tc} (\texttt{mom\_tc}): month * time
    control * opening, 130,672 observations.
  \end{itemize}
\item
  Time-usage tables from \texttt{moves\_meta}:

  \begin{itemize}
  \tightlist
  \item
    \texttt{moves\_month\_tc} (\texttt{mvm\_tc}) month * time control,
    72 observations.
  \item
    \texttt{player\_time\_month\_tc} (\texttt{ptm\_tc}): player * month
    * time control, 1,038,050 observations.
  \end{itemize}
\end{itemize}

The pipeline logs session and version info for both R and Python and
records checksums for downloaded raw files to help with reproducibility.

\subsection{Variables}\label{variables}

This project is working with multiple related-yet-distinct datasets. As
such, it's very difficult to properly introduce the data without an
according quantity of tables, which results in several pages worth of
figures. My personal recommendation is to read through table 1, featured
below, before moving straight to
Section~\ref{sec-observations}.\footnote{I'm considering consolidating
  the Variables and Observations sections into one, as I feel they're
  somewhat redundant, or else perhaps moving the Variable tables
  entirely into the appendix. I need to consult you for permission on
  that first, though.} Most variable names are self-explanatory, so this
section should largely be treated as a reference to skim through if a
variable introduced later is confusing or unclear. Additionally, I
elected to exclude tables \texttt{mm\_tc} and \texttt{mvm\_tc} from this
section, as they're summary-level datasets that won't be explicitly
modeled. Please refer to the \textless appendix placeholder; update when
written\textgreater{} for information on them.

\begin{xltabular}{\textwidth}{>{\raggedright\arraybackslash}p{3.2cm} >{\raggedright\arraybackslash}p{2.8cm} X}
\caption{Core identifiers and categorical variables used across derived datasets.}\label{tab:variables-core}\\
\toprule
\textbf{Variable} & \textbf{Type} & \textbf{Description} \\
\midrule
\endfirsthead

\toprule
\textbf{Variable} & \textbf{Type} & \textbf{Description} \\
\midrule
\endhead

\midrule
\multicolumn{3}{r}{\small\emph{Continued...}} \\
\endfoot

\bottomrule
\endlastfoot

\texttt{player} & categorical & Deterministic pseudonymous user ID (e.g., \texttt{u\us ...}). Appears in player-level tables (\texttt{pm}, \texttt{pm\us tc}, \texttt{pom\us tc}, \texttt{ptm\us tc}). \\
\texttt{month} & categorical & Month of play as \texttt{YYYY-MM}. Primary time index for metagame and longitudinal analysis. \\
\texttt{tc\us bucket} & categorical & Time control bucket: \texttt{blitz}, \texttt{rapid}, or \texttt{classical}. Used to stratify behavior by constraint level. \\
\texttt{eco} & categorical & ECO opening code (e.g., \texttt{C41}). Standardized opening taxonomy useful for grouping openings into families. \\
\texttt{opening} & categorical & Human-readable opening family name (e.g., \texttt{Philidor Defense}). Used in repertoire and metagame summaries. \\
\end{xltabular}


\begin{xltabular}{\textwidth}{>{\raggedright\arraybackslash}p{3.2cm} >{\raggedright\arraybackslash}p{2.6cm} X}
\caption{Player-month outcome and covariate variables (tables: \texttt{pm}, \texttt{pm\us tc}).}\label{tab:variables-player-month}\\
\toprule
\textbf{Variable} & \textbf{Type} & \textbf{Description} \\
\midrule
\endfirsthead

\toprule
\textbf{Variable} & \textbf{Type} & \textbf{Description} \\
\midrule
\endhead

\midrule
\multicolumn{3}{r}{\small\emph{Continued...}} \\
\endfoot

\bottomrule
\endlastfoot

\texttt{games} & numeric (count) & Number of games played in the month (and time-control bucket for \texttt{pm\us tc}). Also a reliability/weighting signal. \\
\texttt{wins} & numeric (count) & Count of wins in the period. \\
\texttt{losses} & numeric (count) & Count of losses in the period. \\
\texttt{draws} & numeric (count) & Count of draws in the period. \\
\texttt{win\us rate} & numeric (proportion) & Key performance metric: wins divided by games. Natural response variable for performance models. \\
\texttt{draw\us rate} & numeric (proportion) & Draws divided by games. Useful because draw propensity varies by skill, time control, and strategic choices. \\
\texttt{avg\us rating} & numeric & Primary skill proxy: average of the player’s pre-game rating across games in the period. Used for stratification and longitudinal change. \\
\texttt{avg\us opp\us rating} & numeric & Average opponent rating. Contextualizes performance and approximates strength of schedule. \\
\texttt{avg\us rating\us diff} & numeric & Average (player rating minus opponent rating). Summarizes typical matchup difficulty. \\
\texttt{avg\us moves} & numeric & Average game length (moves). Proxy for style/complexity. \\
\texttt{opening\us diversity} & numeric (count) & Number of distinct openings used in the period. Proxy for repertoire breadth. \\
\texttt{top\us opening} & categorical & Most frequently played opening during the period (ties broken deterministically). \\
\end{xltabular}

\begin{xltabular}{\textwidth}{>{\raggedright\arraybackslash}p{3.2cm} >{\raggedright\arraybackslash}p{2.6cm} X}
\caption{Player-opening table variables (table: \texttt{pom\us tc}).}\label{tab:variables-player-opening}\\
\toprule
\textbf{Variable} & \textbf{Type} & \textbf{Description} \\
\midrule
\endfirsthead

\toprule
\textbf{Variable} & \textbf{Type} & \textbf{Description} \\
\midrule
\endhead

\midrule
\multicolumn{3}{r}{\small\emph{Continued...}} \\
\endfoot

\bottomrule
\endlastfoot

\texttt{games} & numeric (count) & Number of games the player played with the given opening in the month/time-control bucket. Used as a frequency weight. \\
\texttt{win\us rate} & numeric (proportion) & Opening-specific win rate for the player and period. Used to study effectiveness conditional on opening. \\
\texttt{draw\us rate} & numeric (proportion) & Opening-specific draw rate. Helps characterize risk profile and style. \\
\texttt{avg\us rating} & numeric & Average player rating in games where the opening was played. Enables comparison of opening choice across skill. \\
\texttt{avg\us rating\us diff} & numeric & Average rating difference in games where the opening was played. Captures whether openings are used in tougher/easier matchups. \\
\end{xltabular}


\begin{xltabular}{\textwidth}{>{\raggedright\arraybackslash}p{3.2cm} >{\raggedright\arraybackslash}p{2.6cm} X}
\caption{Player time-use table variables (table: \texttt{ptm\us tc}).}\label{tab:variables-player-time}\\
\toprule
\textbf{Variable} & \textbf{Type} & \textbf{Description} \\
\midrule
\endfirsthead

\toprule
\textbf{Variable} & \textbf{Type} & \textbf{Description} \\
\midrule
\endhead

\midrule
\multicolumn{3}{r}{\small\emph{Continued...}} \\
\endfoot

\bottomrule
\endlastfoot

\texttt{move\us rows} & numeric (count) & Number of move observations contributing to the summary (after sampling and filtering). Reliability signal. \\
\texttt{games} & numeric (count) & Number of distinct games contributing move records for the player-month-time-control group. \\
\texttt{avg\us time\us spent} & numeric (seconds) & Mean time spent per move. Captures typical deliberation time. \\
\texttt{median\us time\us spent} & numeric (seconds) & Median time spent per move; robust to outliers and occasional long thinks. \\
\texttt{p90\us time\us spent} & numeric (seconds) & 90th percentile of time spent per move. Measures tendency for occasional deep computation. \\
\texttt{avg\us clock\us after} & numeric (seconds) & Mean remaining clock time after moves. Proxy for time management and pacing. \\
\texttt{low\us time\us rate} & numeric (proportion) & Fraction of moves under a low-time threshold. Proxy for time trouble frequency. \\
\texttt{avg\us rating} & numeric & Average player rating across move records in the group; allows stratification by skill. \\
\texttt{avg\us rating\us diff} & numeric & Average rating difference across the move sample; controls for opponent strength effects. \\
\texttt{avg\us time\us opening} & numeric (seconds) & Average time spent per move during the opening phase (if phase segmentation is available; game phases are defined by chess-engine analysis). \\
\texttt{avg\us time\us middlegame} & numeric (seconds) & Average time spent per move during middlegame phase (may be missing for very short games). \\
\texttt{avg\us time\us endgame} & numeric (seconds) & Average time spent per move during endgame phase (often missing because many games never reach endgame). \\
\end{xltabular}

\begin{xltabular}{\textwidth}{>{\raggedright\arraybackslash}p{3.2cm} >{\raggedright\arraybackslash}p{2.6cm} X}
\caption{Metagame opening table variables (table: \texttt{mom\us tc}).}\label{tab:variables-metagame-opening}\\
\toprule
\textbf{Variable} & \textbf{Type} & \textbf{Description} \\
\midrule
\endfirsthead

\toprule
\textbf{Variable} & \textbf{Type} & \textbf{Description} \\
\midrule
\endhead

\midrule
\multicolumn{3}{r}{\small\emph{Continued...}} \\
\endfoot

\bottomrule
\endlastfoot

\texttt{games} & numeric (count) & Number of sampled games for the opening within month and time control. Precision/support indicator. \\
\texttt{share} & numeric (proportion) & Core metagame measure: fraction of games using the opening within a month/time-control bucket. Tracks adoption shifts. \\
\texttt{white\us win\us rate} & numeric (proportion) & Fraction of games where White wins. Used for opening-level success patterns and color asymmetries. \\
\texttt{black\us win\us rate} & numeric (proportion) & Fraction of games where Black wins. Complements White win rate. \\
\texttt{draw\us rate} & numeric (proportion) & Fraction of games drawn. Draw propensity varies by opening and time control. \\
\texttt{upset\us rate} & numeric (proportion) & Fraction of games where the lower-rated player wins (pipeline definition). Proxy for volatility/competitiveness. \\
\texttt{avg\us abs\us rating\us diff} & numeric & Mean absolute rating gap between players in the opening/month/tc group. Context for outcome rates. \\
\texttt{avg\us moves} & numeric & Mean game length for the opening/month/tc group. Proxy for complexity and style at population level. \\
\end{xltabular}

\subsection{Observations}\label{sec-observations}

Each derived dataset is organized at a different unit of analysis (e.g.,
player-by-month or opening-by-month). In all cases, the raw unit of data
is a rated chess game from the Lichess public database, but the pipeline
aggregates games (and, for time-use, moves) into analysis-ready
summaries. Because many variables overlap across tables, the observation
structure can be summarized with a small set of representative row
types. The examples below demonstrate the core units: (1) player-period
summaries (with optional time-use extensions), (2) player opening-level
summaries, and (3) population-level metagame summaries for openings.

\begin{xltabular}{\textwidth}{>{\raggedright\arraybackslash}p{3.4cm} X}
\caption{Sample observation for player-period outcome summaries (tables: \texttt{pm}, \texttt{pm\us tc}).}\label{tab:obs-player-period-outcomes}\\
\toprule
\textbf{Field} & \textbf{Example Value/Meaning} \\
\midrule
\endfirsthead

\toprule
\textbf{Field} & \textbf{Example Value/Meaning} \\
\midrule
\endhead

\midrule
\multicolumn{2}{r}{\small\emph{Continued...}} \\
\endfoot

\bottomrule
\endlastfoot

\texttt{player} & \texttt{u\us 0001cf80083d0517} (stable pseudonymous identifier) \\
\texttt{month} & \texttt{2024-12} (calendar month) \\
\texttt{tc\us bucket} & \texttt{rapid} (present in \texttt{pm\us tc}; omitted in \texttt{pm}) \\
\texttt{games} & 5 (total games in this player-month-(tc) group) \\
\texttt{wins} & 4 \\
\texttt{losses} & 0 \\
\texttt{draws} & 1 \\
\texttt{win\us rate} & 0.80 (wins / games) \\
\texttt{draw\us rate} & 0.20 (draws / games) \\
\texttt{avg\us rating} & 1544.0 (mean pre-game rating across games) \\
\texttt{avg\us opp\us rating} & 1503.2 (mean opponent rating across games) \\
\texttt{avg\us rating\us diff} & 40.8 (mean rating difference: player minus opponent) \\
\texttt{avg\us moves} & 40.4 (mean game length in moves) \\
\texttt{opening\us diversity} & 4 (distinct openings played in the group) \\
\texttt{top\us opening} & \texttt{Ponziani Opening: Jaenisch Counterattack} (most frequent opening) \\
\end{xltabular}

\begin{xltabular}{\textwidth}{>{\raggedright\arraybackslash}p{3.4cm} X}
\caption{Sample observation for player-period time-use summaries (table: \texttt{ptm\us tc}).}\label{tab:obs-player-timeuse}\\
\toprule
\textbf{Field} & \textbf{Example Value/Meaning} \\
\midrule
\endfirsthead

\toprule
\textbf{Field} & \textbf{Example Value/Meaning} \\
\midrule
\endhead

\midrule
\multicolumn{2}{r}{\small\emph{Continued...}} \\
\endfoot

\bottomrule
\endlastfoot

\texttt{player} & \texttt{u\us 00004578077e9fa4} (stable pseudonymous identifier) \\
\texttt{month} & \texttt{2024-08} \\
\texttt{tc\us bucket} & \texttt{blitz} \\
\texttt{move\us rows} & 70 (number of move records contributing to this summary) \\
\texttt{games} & 1 (distinct games represented among the sampled moves) \\
\texttt{avg\us time\us spent} & 1.66 seconds (mean time spent per move) \\
\texttt{median\us time\us spent} & 1.0 seconds (median time spent per move) \\
\texttt{p90\us time\us spent} & 3.1 seconds (90th percentile time spent per move) \\
\texttt{avg\us clock\us after} & 122.5 seconds (mean remaining clock after moves) \\
\texttt{low\us time\us rate} & 0.00 (share of moves flagged as low-time) \\
\texttt{avg\us rating} & 1677.0 (mean pre-move rating across sampled moves) \\
\texttt{avg\us rating\us diff} & 11.0 (mean rating difference across sampled moves) \\
\texttt{avg\us time\us opening} & 0.9 seconds (mean time spent in opening phase) \\
\texttt{avg\us time\us middlegame} & 2.2 seconds (mean time spent in middlegame phase) \\
\texttt{avg\us time\us endgame} & 1.64 seconds (mean time spent in endgame phase; may be missing) \\
\end{xltabular}

\begin{xltabular}{\textwidth}{>{\raggedright\arraybackslash}p{3.4cm} X}
\caption{Sample observations for opening-level summaries: player repertoire entries (\texttt{pom\us tc}) and population metagame cells (\texttt{mom\us tc}).}\label{tab:obs-opening-level}\\
\toprule
\textbf{Field} & \textbf{Example Value/Meaning} \\
\midrule
\endfirsthead

\toprule
\textbf{Field} & \textbf{Example Value/Meaning} \\
\midrule
\endhead

\midrule
\multicolumn{2}{r}{\small\emph{Continued...}} \\
\endfoot

\bottomrule
\endlastfoot

\multicolumn{2}{l}{\textbf{Player opening summary (table: \texttt{pom\us tc})}} \\
\texttt{player} & \texttt{u\us 0001cf80083d0517} \\
\texttt{month} & \texttt{2024-11} \\
\texttt{tc\us bucket} & \texttt{blitz} \\
\texttt{eco} & \texttt{B13} \\
\texttt{opening} & \texttt{Caro-Kann Defense: Exchange Variation} \\
\texttt{games} & 4 (games where this player used the opening in this month/tc) \\
\texttt{win\us rate} & 0.50 (opening-conditional win rate for the player) \\
\texttt{draw\us rate} & 0.00 (opening-conditional draw rate for the player) \\
\texttt{avg\us rating} & 1965.3 (mean rating in games using this opening) \\
\texttt{avg\us rating\us diff} & 12.0 (mean rating difference in those games) \\

\midrule
\multicolumn{2}{l}{\textbf{Population opening summary (table: \texttt{mom\us tc})}} \\
\texttt{month} & \texttt{2024-01} \\
\texttt{tc\us bucket} & \texttt{blitz} \\
\texttt{eco} & \texttt{C41} \\
\texttt{opening} & \texttt{Philidor Defense} \\
\texttt{games} & 8{,}825 (sampled games using this opening in the month/tc group) \\
\texttt{share} & 0.0189 (opening share of the month/tc metagame) \\
\texttt{white\us win\us rate} & 0.512 (fraction of games White wins) \\
\texttt{black\us win\us rate} & 0.442 (fraction of games Black wins) \\
\texttt{draw\us rate} & 0.0466 (fraction drawn) \\
\texttt{upset\us rate} & 0.436 (fraction where lower-rated player wins) \\
\texttt{avg\us abs\us rating\us diff} & 50.2 (mean rating gap magnitude) \\
\texttt{avg\us moves} & 35.5 (mean game length) \\
\end{xltabular}

Together, these observation types support longitudinal player
performance modeling (\texttt{pm}, \texttt{pm\_tc}), repertoire/style
measurement from opening choice (\texttt{pom\_tc}), metagame dynamics
via opening adoption and outcomes (\texttt{mom\_tc}), and
bounded-rationality/time-management analyses using move-time features
(\texttt{ptm\_tc}).

\subsection{Cleaning}\label{cleaning}

A large amount of cleaning was performed during the pipeline's parsing
steps. The filtering and deterministic sampling was already touched upon
in \#sec-source, but additional standardization and feature construction
include:

\begin{itemize}
\tightlist
\item
  \texttt{TimeControl} tag parsed into:

  \begin{itemize}
  \tightlist
  \item
    \texttt{initial} (seconds).
  \item
    \texttt{increment} (seconds).
  \item
    \texttt{tc\_bucket} (\textless{} 180 = bullet, \textless{} 600 =
    blitz, \textless{} 1800 = rapid, else classical),
  \end{itemize}
\item
  Converted PGN result strings into:

  \begin{itemize}
  \tightlist
  \item
    \texttt{winner} (white/black/draw).
  \item
    \texttt{score\_white} (1/0.5/0).
  \end{itemize}
\item
  Computed upset indicators using ratings:

  \begin{itemize}
  \tightlist
  \item
    \texttt{rating\_diff} = \texttt{white\_rating} -
    \texttt{black\ rating}.
  \item
    \texttt{abs\_rating\_diff}
  \item
    \texttt{is\_upset}
  \end{itemize}
\item
  Game-length variables:

  \begin{itemize}
  \tightlist
  \item
    Computed \texttt{piles} and \texttt{moves} from move sequences (used
    to construct final datasets).
  \end{itemize}
\end{itemize}

The information was then restructured and aggregated into a tidy format.

\section{Exploratory Data Analysis}\label{exploratory-data-analysis}

Show the initial graphs and basic analysis methods used to determine
your methodology

\subsection{Interesting EDA 1 (Change
name)}\label{interesting-eda-1-change-name}

Show some graph or numerical summary that imparts an interesting idea
about the data set

\subsection{Interesting EDA 2 (Change
name)}\label{interesting-eda-2-change-name}

Show some graph or numerical summary that imparts an interesting idea
about the data set

\newpage

\section{Methodology}\label{methodology}

Describe the methods used to analyze the data. Do not discuss the
analysis itself. Make sure to discuss any train/test data splitting
here.

\subsection{Data Tranformations (if
applicable)}\label{data-tranformations-if-applicable}

Describe any transformation you did to the data to assist in the
analyzing of the data such as creating dummy variables.

\subsection{Types of Models}\label{types-of-models}

Describe how many different models you used and why.

\subsubsection{Model 1 (Change to appropriate
name)}\label{model-1-change-to-appropriate-name}

Describe why you used model type one. You do not need to include a new
section for every iteration of the same model type. For example you do
not need 4 sections if you have 4 linear models with 1, 2, 3, and 4
independent variables.

\subsubsection{Model 2 (Change to appropriate
name)}\label{model-2-change-to-appropriate-name}

Describe why you used model type two.

\subsubsection{Model 3 (Change to appropriate
name)}\label{model-3-change-to-appropriate-name}

Describe why you used model type three.

\subsubsection{Model 4 (Change to appropriate
name)}\label{model-4-change-to-appropriate-name}

Describe why you used model type four.

\subsubsection{Model 5 (Change to appropriate
name)}\label{model-5-change-to-appropriate-name}

Describe why you used model type five.

\subsubsection{Further Models}\label{further-models}

If using more than 5 types of models, use this section to quickly
discuss them here. These should only be models that you do not expect to
be great.

\subsection{Method of model selection}\label{method-of-model-selection}

Describe how you determine which model you choose at the end. Should
also describe why you chose this method.

\newpage

\section{Results}\label{results}

Describe the results that you received from the methods described
earlier.

\subsection{Model 1 (Change to appropriate
name)}\label{model-1-change-to-appropriate-name-1}

Show the results for model type one. At least one fully written out
model should be displayed for each model type.

\subsection{Model 2 (Change to appropriate
name)}\label{model-2-change-to-appropriate-name-1}

Show the results for model type two.

\subsection{etc}\label{etc}

\subsection{Selected model}\label{selected-model}

Display the table that shows the model selection. If using more than 10
models only show the top 10 performing in the table, the full table
should be in the appendix. The final model should be fully displayed.

\newpage

\section{Discussion}\label{discussion}

\subsection{Final model interpolation}\label{final-model-interpolation}

Discuss what information the selected model provides that may be new and
interesting

\subsection{Use of Model}\label{use-of-model}

Describe how this model might be used by others.

\subsection{Comparison}\label{comparison}

Compare the things learned from this selected model to the previous
information.

\newpage

\section{Furture Work}\label{furture-work}

Discuss what work could be done with this model to make it better.

\newpage

\section{References}\label{references}

List the reference used for your paper. This should include anything
cited in your background section and the source of your data.

\newpage

\section{Appendix}\label{appendix}

\subsection{Appendix A Graphs (Use appropriate
name)}\label{appendix-a-graphs-use-appropriate-name}

\subsubsection{A.1 (Use appropriate
name)}\label{a.1-use-appropriate-name}

If showing multiple graphs make sure to index them by 0.1 each

\newpage

\subsection{Appendix B Tables (Use appropriate
name)}\label{appendix-b-tables-use-appropriate-name}

\subsubsection{Appendix B.1 (Use appropriate
name)}\label{appendix-b.1-use-appropriate-name}

If showing multiple tables make sure to index them by 0.1 each

\subsection{Code}\label{code}

\begin{Shaded}
\begin{Highlighting}[]
\FunctionTok{library}\NormalTok{(tidyverse)}
\CommentTok{\# Reason for library ...}
\CommentTok{\#library(...)}
\CommentTok{\# Reason for library ...}
\CommentTok{\#library(...)}

\CommentTok{\# Set seed is used because question ... uses a random generator}
\FunctionTok{set.seed}\NormalTok{(}\DecValTok{123}\NormalTok{)}
\CommentTok{\# Load cleaned data here. If cleaning was required make sure to have a separate r file for the cleaning.}
\CommentTok{\# Use this code to display useful graphs that do not belong in the body of the paper.}
\CommentTok{\# For example use this area for model validation graphs like residual versus fit.}
\CommentTok{\# Use this area to show results that are in table format}
\CommentTok{\# Write your cleaning code here}
\end{Highlighting}
\end{Shaded}



  \bibliography{references.bib}



\end{document}
